% critique.tex — f175r, vv. 1045–1081 : commentaire critique.
% Parties textuelles du commentaire de l'auteur (travaux antérieurs, Harmattan 2023), texte et notes
% conservés ; répartition, harmonisation et passages signalés : repartition.md.
\cc{1054–5}{λαιμότμητ᾽ ἄχη δαιμόνων αἰόλαν / νέκυσιν ἐπὶ πορείαν.}{Willems a consacré une note à ce passage\nc{\nieme{27} note aux \italique{Thesmophories} d'Alphonse Willems, reproduite en annexe.}. Si les vues qu'il y développe sont ingénieuses, elle n'emporte pas pour autant ma conviction. Willems propose de construire ἄχη δία μένων là où le \italique{Ravennas} porte ἄχη †δαιμόνων†. À défaut d'autres propositions plus convaincantes, c'est néanmoins la correction que je retiendrai.}
\cc{1059}{Ἠχὼ λόγων ἀντῳδὸς ἐπικοκκάστρια.}{Pour le sens d'ἐπικοκκάστρια, voir le commentaire de la pièce.\\
Il est étonnant que les modernes éditions n'aient pas cherché à percer ce que les répliques d'Écho-Euripide pouvaient avoir d'amusantes ou de méchantes… Pourtant, elles disposaient déjà d'un exemple de détournement par la « rediseuse goguenarde » : à la fin de la scène, au vers 1094, alors que l'archer Scythe tente d'attraper Écho-Euripide et que ce dernier s'enfuit sans demander son reste, une scholie nous indiquait clairement ce que disait en son jargon l'archer Scythe :\citComm{ἀντὶ τοῦ `οὐ χαιρήσεις'\\
À la place de : « tu ne vas pas te réjouir ».}On peut déduire de cette scholie que le Scythe – qui ne tient aucun compte dans son parler des aspirations – avait menacé Écho-Euripide d'un : `οὐ καιρήσεις'. Or, le manuscrit de Ravenne comportait effectivement la réponse ἐπικοκκάστρια d'Euripide-Écho\nc{En attribuant, il est vrai, la réplique à l'archer. Mais les indications de personnages sont de peu de valeur : les anciens parchemins ne comportaient pour la plupart que de légers signes (\italique{paragraphos} ou parfois simple indentation) pour indiquer le changement d'interlocuteur, et ce sans préciser quels étaient ces interlocuteurs ; les indications de personnages furent ajoutées par des scribes très postérieurement… et sont parfois assez folkloriques. (Pour celles du \italique{Ravennas}, quelques indications sont de la main d'Euphrosynus Boninus, celui qui compulsa le manuscrit pour la première édition imprimée des \italique{Thesmophories} -- la \italique{Juntina} (Venise, 1516) des apparats critiques. Von Velsen sut identifier les annotations de la main d'Euphrosynus Boninus sur le \italique{Ravennas} et ne les fit pas apparaître dans son apparat critique. Je l'ai suivi en cela. Cf. Albert Martin, \italique{Les scolies du manuscrit d'Aristophane à Ravenne} (1882), p. V ; A. von Velsen, \italique{Über den Codex Urbinas der Lysistrata und der Thesmophoriazusen des Aristophanes} (1871).)} : `οὐκ᾽ αἰρήσεις', c'est-à-dire en bon attique : `οὐχ αἱρήσεις', « tu ne m'attraperas pas ! ». Le cas était simple mais il aura fallu tout de même des siècles de philologiques hésitations entre les répartitions : Τοξότης (l'Archer) : « οὐ καιρήσεις » — Εὐριπίδης : « οὐ καιρήσεις » et Τοξότης : « οὐκ αἰρήσεις » — Εὐριπίδης : « οὐκ αἰρήσεις » avant que le grand Wilamowitz\nc{En simple note de ses \italique{Kleine Schriften}, IV (1962), p. 486. Mais ces écrits n'ont été publiés qu'après la mort du grand philologue. Si une partie des travaux modernes rétablissent correctement la menace du Scythe et son écho, il semble que le mérite en revienne plutôt à Sommerstein, qui eut la bonne intuition indépendamment de Wilamowitz.} ne s'avise de la solution conciliatrice, la seule sensée et la seule comique : Τοξότης : « οὐ καιρήσεις » — Εὐριπίδης : « οὐχ αἱρήσεις ».\\
C'est -- reconnaissons-le -- un exercice des plus périlleux que de s'aventurer dans des tentatives de reconstruction ; laisser ainsi libre cours à l'imagination peut exposer à de rudes retours de bâton ; pour Aristophane, nous devons cependant impérieusement nous y livrer puisqu'il vient de nous prévenir que les échos seraient \italique{moqueurs} ! Avec une réserve cependant : lorsque l'ἐπικοκκάστρια d'une réplique m'est demeurée celée, et que malgré mes efforts les seules propositions que j'avais à présenter n'étaient pas même dignes d'intégrer l'apparat critique -- le cas n'est pas isolé --, je lui ai apposé la \italique{croix du désespoir} †.. Charge à d'autres de compléter mon ébauche ! je pense, pour ma part, ne pas avoir démérité d'Aristophane.}
\cc{1069}{τοῦ σεμνοτάτου δι᾽ Ὀλύμπου; --- δι᾽ Ὀλύμπου.}{En bonne rigueur philologique, je dois me demander sur quels éléments probants s'appuyer dans mon aventureuse entreprise de rétablissement de l'\allemand{Urtext} des \italique{Thesmophories}. Car il ne suffit pas de partir à l'aventure, toutes brides de l'imagination abattues ; encore faut-il apporter des indices de ce que j'avance. 
J'ai déjà évoqué des bribes de scholies ; nous en croiserons d'autres en route. 
Une variante orthographique entre un mot et son écho nous fournira aussi une amorce philologique. 
Cela est bien maigre... est-ce qu'Hésychius, ou les commentateurs qu'il recopie\nc{Cf. le passionnant livre posthume de Pierre Boudreaux, \italique{Le texte d'Aristophane et ses commentateurs}.}, dont j'ai assez montré toute l'importance, avait entre les mains une édition d'Aristophane assez bien conservée pour qu'y apparaissent certains détournements de la nymphe Écho ? Réponse : ce n'est pas impossible ! voyons cela avec la première réplique ἐπικοκκάστρια de la nymphe Écho. \\
Voici la glose, parcellaire, d'Hésychius (ο 649) qui éclaire et justifie peut-être l'interruption d'Écho :\citComm{ὀλυμπία· δοκός\\
olympia : poutre.}La langue d'Aristophane possède la matrice métaphorique : \italique{poutre} $=$ \italique{poutre à laquelle on se pend} $=$ \italique{pendaison}. (Le français dispose d'une matrice semblable avec le mot \italique{corde}.) J'en veux pour preuves : \ \ \ 1° \ ce commentaire\nc{Fragment d'Aristophane 528 K.-A. (515 Kock). Cf. l'édition Kassel-Austin des \italique{Poetae Comici Graeci} pour la juste distinction entre δοκός oxyton et δόκος paroxyton.} : [δοκὸς] παρὰ Ἀριστοφάνει ἐν Ταγηνισταῖς σημαίνει τὴν ἀγχόνην, commentaire repris par Hésychius (δ 2129) : δοκόν· τὴν ἀγχόνην.
\ \ \ 2° \ Dans les \italique{Oiseaux}, après avoir poutré l'importun chresmologue Hiéroclès, on l'envoie se faire pendre εἰς Ἐλύμνιον (vers 1126), ville d'Eubée dont le nom fait entendre à l'oreille grecque sa \italique{poutre} : ἐλύμνιαι· δοκοὶ ὀρόφιναι (Hésychius δ 2227).\\
Rappeler à Andromède sa condition de crucifiée \italique{pendue} à la falaise, voilà une incise qui pourrait être qualifiée de κακοστένακτος.\\
Je m'arrête ici, au bord de la reconstruction tentatrice. Car grande est la tentation de retoucher au texte même d'Andromède en y intégrant le iota d'ὀλυμπία et voir dans `δι᾽ Ὀλύμπου' une mention de Ζεύς Ὀλύμπιος\nc{Διός Ὀλυμπίου \italique{contra metrum} cependant.} ; qu'on se souvienne du vers 272 : Euripide jurait sur l'\italique{Éther, demeure de Zeus} (αἰθέρ᾽, οἴκησιν Διός\nc{Aristophane y revient au vers 100 des \italique{Grenouilles} : αἰθέρα Διὸς δωμάτιον.}). Le vrai Euripide parle ailleurs du \italique{ciel étoilé de Zeus} (\italique{Ion}, v. 1077 sq.) :\citComm{καὶ Διὸς ἀστερωπὸς\\
ἀνεχόρευσεν αἰθήρ.}La mention de Zeus lorsqu'un poète évoque le ciel constellé est d'ailleurs très naturelle puisque le quatrième astre le plus lumineux, après le Soleil, la lune et Vénus, n'est autre que \italique{Jupiter}, Διὸς ἀστήρ.}
\cc{1071}{μέρος ἐξέλαχον;}{Sans \italique{tordre} le texte, il me semble que la réplique ἐπικοκκάστρια d'Écho s'entend assez bien -- soit dit avec toute la prudence qu'une telle chirurgie du texte impose :\citComm{μέρος ἐξέλαχον --- μόρος ἔξελ(ε)\nc{Aoriste \italique{gnomique}.} ἄχος\\
La Mort\nc{Cette acception de μόρος a été relevée par à peu près tous les lexicographes anciens : Hésychius μ 1683 ; Photius μ 538 ; Συναγωγὴ λέξεων χρησίμων μ 267 ; Suidas μ 1257 ; \latin{Etymologicum genuinum} \italique{sub voce} μόρος· ὁ θάνατος ; \&c.} extirpe la douleur.}Pour le sens et les parallèles, voir le commentaire de la pièce.}
\cc{1073–5}{ὦ γραῦ στωμυλλωμένη — στωμυλωμένη — νὴ Δί᾽.}{Le manuscrit de Ravenne (qui, rappelons-le, est le seul témoin complet de notre pièce) comporte deux orthographes différentes pour στωμυλλωμένη : avec un double lambda -λλ- au vers 1073 et un unique -λ- au vers 1074... Étourderie ou fantaisie de notre scribe jonglant avec une orthographe dont il est peu sûr ? Possible... ou bien preuve qu'il recopiait scrupuleusement un texte sans parfois – comme ici – le comprendre ? c'est ce que je vais tenter de montrer : il se trouve que le verbe στωμύλλω, \italique{babiller}, \italique{déblatérer}, s'écrit effectivement avec deux lambdas en grec classique. Cependant, l'adjectif στωμύλος (\italique{qui parle beaucoup}, et souvent avec nuance méliorative par rapport au verbe) qui en dérive n'en comporte qu'un seul ! Si donc, l'on fait le pari que le scribe ne s'y est pas trompé en conservant l'orthographe, savoir : on est en présence d'un participe moyen-passif au vers 1073 et de l'adjectif au vers 1074, il est alors aisé de retrouver la véritable réponse \italique{ἐπικοκκάστρια} d'Écho : στωμυλιο-μενής, « [je suis] de tempérament prolixe ».
Le composé est de la même facture que στωμυλιο-συλλεκτάδης, hapax qu'Aristophane mit dans la bouche d'Eschyle au vers 839 des \italique{Grenouilles} pour clore le bec à Euripide, cet intarissable \italique{collectionneur de commérages}. La seconde partie du composé est forgée avec le suffixe -μενής, \italique{au tempérament}, \italique{à la disposition...} (sur le modèle d'εὐ-μενής, δυσ-μενής, πρευ-μενής \&c).\\
Ma conjecture est renforcée par la concessive de Mnésiloque au vers suivant : il abonde, νὴ Δί᾽, ... néanmoins, γ᾽ : « ça oui par Zeus ! du moins... ».}
\cc{1078}{καὶ χαριεῖ μοι παῦσαι --- παῦσαι.}{Mnésiloque-Andromède, qui ne peut chanter sa monodie correctement, demande à Euripide-Écho de cesser de l'interrompre :\citComm{καὶ χαριεῖ μοι παῦσαι...}«~Tu m'obligeras d'arrêter...~» -- «~De crier~» interrompt facilement Écho. (Et, bien-sûr, en le lui criant à l'oreille :)\citComm{ (ἐ)παΰσαι.}(Le verbe ἐπαΰω, attesté par Eschyle dans ses \italique{Choéphores} (v. 828), est forgé avec le préfixe ἐπι- à l'instar d'ἐπικοκκάστρια et le verbe αὔω, \italique{crier} (cf. le vers 1019).)\\
Une telle provocation explique la réaction d'emportement de Mnésiloque :\citComm{βάλλ᾽ ἐς κόρακας.}}
\cc{1079}{βάλλ᾽ ἐς κόρακας — βάλλ᾽ ἐς κόρακας.}{Dans mon entreprise fort périlleuse de dénicher des jeux de mots tout au long de ce passage, en voici un qui est moins spéculatif et qui va nous permettre de reprendre haleine ! Les Grecs usaient d'une formule d'imprécation : ἐς κόρακας, équivalente de notre « [va] au diable ! ». Aristophane l'emploie fréquemment avec le verbe βάλλειν : \italique{Guêpes} v. 835 ; \italique{Ploutos} v. 782 ; \italique{Nuées} v. 133 – digressons d'ailleurs sur ce passage des \italique{Nuées} puisqu'Aristophane recourt au même jeu de mots qu'ici : Strepsiade arrive devant le pensoir de Socrate et toque à la porte, il se fait recevoir par un peu accueillant :\citComm{Βάλλ᾽ ἐς κόρακας· / Τίς ἐσθ᾽ ὁ κόψας τὴν θύραν;}signalons à toutes fins utiles que les κόρακες dont parle le disciple de Socrate sont des marteaux de porte\nc{Sens attesté par Pollux, X, 22.20.}. 
Le disciple tint donc à peu près ce langage : \citComm{Βάλλε κόρακας·/ Τίς ἐσθ᾽ ὁ κόψας τὴν θύραν;\\
Bas les marteaux ! Qui est ce tambourineur de porte ?}(compatible avec le mètre iambique et surtout avec le bon sens : qui – fût-il disciple de Socrate – s'aviserait de demander l'identité de quelqu'un qu'on vient d'envoyer se faire voir chez les corbeaux grecs ?)\\
Il en va de même dans le passage qui nous intéresse présentement : Écho conseille frivolement à Mnésiloque de retirer les crochets qui le retiennent : « Mais va te faire voir ! — Et toi mets bas le fermoir ! »}
\cc{1080}{τί τὸ κακόν;}{La leçon du \italique{Ravennas}, «~τί τὸ κακόν; --- τί τὸ κακόν;~», ne respecte pas le mètre anapestique. Gardons-nous pour autant de supprimer avec Bentley l'article `τὸ'. 
Si la formule «~τί τὸ κακόν;~» (ou d'autres très approchantes) se rencontre couramment chez Aristophane\nc{\italique{Thesmophories}, v. 610 ; \italique{Paix}, v. 322 ; \italique{Oiseaux}, v. 1213 ; \italique{Guêpes} v. 1136, τουτὶ τὸ κακὸν τί ἐστι ; \italique{Paix}, v. 181.}, sa version lapidaire, «~τί κακόν;~», elle, ne se trouve jamais\nc{Certains éditeurs prétendent que si : au vers 1085. En vérité, ces faussaires commettent là une double forfaiture en «~corrigeant~» le présent vers \italique{metri gratia} ainsi que le vers 1085 ; le second forfait servant d'alibi au premier : au vers 1085, ils font apparaître «~une seconde fois~» le tour `τί κακόν' alors que le manuscrit porte `σί κακόν', confirmé par le scholiaste (ἀντὶ τοῦ `σοὶ κακόν'). Et quand bien même : c'est dans la bouche d'un Scythe qu'Aristophane la met au vers 1085, l'argument pour la faire dire aussi à Mnésiloque est irrecevable. En dernière analyse, ces éditeurs manquent cruellement de bon sens : le Scythe en son jargon ne tient déjà aucun compte des aspirations, de la grammaire, il ignore sa conjugaison et ses déclinaisons, il iotacise à l'envi... et l'on voudrait encore qu'il sigmatise ses `τ' ?! un tel massacre du grec aurait achevé d'être tout-à-fait inintelligible aux spectateurs.}. La bonne solution n'est donc pas celle de Bentley, et pour ma part je me résigne à \latin{crux}-désespérer †}
